% Journal Paper: Unified Multimodal Graph-Gated MoE for Mental Health Assessment
% Generated: 2026-07-14

\documentclass[twocolumn]{article}
\usepackage[utf8]{inputenc}
\usepackage[T1]{fontenc}
\usepackage{graphicx}
\usepackage{amsmath}
\usepackage{amssymb}
\usepackage{booktabs}
\usepackage{hyperref}
\usepackage{cleveref}
\usepackage{siunitx}
\usepackage{xcolor}
\usepackage{url}
\usepackage[numbers]{natbib}
\usepackage{microtype}
\usepackage{balance}

\graphicspath{{../}}
\DeclareGraphicsExtensions{.png,.jpg,.pdf}

\hypersetup{
    colorlinks=true,
    linkcolor=blue,
    citecolor=blue,
    urlcolor=blue,
}

\title{Graph-Guided Mixture-of-Experts for Unified Multimodal Mental Health Assessment}

\author{
    \textbf{Anilson Silva}\\
    Department of Computer Science\\
    \texttt{anilson.silva@example.com}
}

\begin{document}

\maketitle

\begin{abstract}
We propose a unified multimodal architecture for mental health assessment that jointly learns from depression, sentiment, emotion, and personality across three public benchmarks: DAIC-WOZ (clinical interviews, $n=189$), CMU-MOSEI (YouTube videos, $n=22{,}777$), and ChaLearn First Impressions (video clips, $n=10{,}000$). The model combines modality-specific encoders (RoBERTa for text, WavLM for audio, ViT for video), gated late fusion, a Mixture-of-Experts layer, and a KNN-graph-based GraphSAGE router that selects experts using neighborhood information. Our central finding is that graph routing consistently improves over non-graph baselines: variant V0 achieves MOSEI sentiment CCC=0.6803 (+0.18 over MMoEEx) and variant V3 achieves DAIC depression AUROC=0.8967 (+0.40 over MMoEEx). Cross-attention fusion does not reproduce a reported gain over gated fusion under our evaluation protocol. LLM encoders (Mistral-7B, CLAP, LLaVA) improve over classical encoders but do not surpass graph routing. We release a leakage-safe graph construction protocol, a full ablation ladder, and graph-based explanations for clinical interpretability.
\end{abstract}

\section{Introduction}
\label{sec:introduction}

Mental health assessment from multimodal signals is a cornerstone challenge in affective computing and clinical informatics. Depression, the primary focus of this work, affects over 280 million people worldwide \citep{who2023depression} and is routinely assessed through clinical interviews that combine speech, facial expressions, and verbal content. Meanwhile, related affective signals---sentiment, emotion, and apparent personality---provide complementary supervision signals that can enrich mental health models through multitask learning \citep{dupont2020multitask,kendall2018multitask}.

Our main claim is that a single multimodal architecture can learn both shared and task-specific representations for depression, affect, and personality, and that graph-based routing improves prediction and makes model behavior easier to explain. We do not claim that apparent personality and clinical depression measure the same construct. The proposed system is intended as a decision-support tool and should not replace clinical diagnosis; all results require prospective validation before clinical deployment.

This work makes three contributions:
\begin{enumerate}
    \item \textbf{Graph-gated MoE routing}: Combining KNN graph construction with GraphSAGE-based expert selection, demonstrating that topology-aware expert routing consistently improves over non-graph baselines across three affective computing benchmarks.
    \item \textbf{Leakage-safe graph construction protocol}: Three graph construction modes (inductive, split-local, transductive) under strict subject-independent splits, with formal validation that prevents test-label leakage through graph edges.
    \item \textbf{Comprehensive empirical analysis}: A systematic ablation ladder spanning unimodal baselines, three fusion strategies, MoE variants, five graph configurations, and LLM encoder replacements.
\end{enumerate}

\section{Related Work}
\label{sec:related}

\paragraph{Multimodal Fusion.} Early fusion for mental health used simple concatenation \citep{dupont2020multitask}. Later approaches adopted late fusion with learned modality gates \citep{hazdar2024gated} or Low-Rank Multimodal Fusion (LMF) \citep{zadeh2017LMF}. Cross-modal attention mechanisms have been proposed for depression detection \citep{kim2024crossmodal}, but in our experimental configuration, we do not observe the reported +0.041 AUROC gain.

\paragraph{Mixture of Experts.} The sparsely-gated MoE layer \citep{shazeer2017moe} routes each input to a subset of experts. Multitask MoE (MMoE) extends this with task-specific gates \citep{ma2018model}, and MMoEEx adds expert exclusivity and orthogonality regularization \citep{jacobs2024mmoe}. We build on MMoEEx but add graph-based routing.

\paragraph{Graph Neural Networks for Routing.} GraphSAGE \citep{graphsage2017inductive} performs inductive representation learning on graphs. We use GraphSAGE as a router: aggregating features from K nearest neighbors to produce a graph context vector that modulates expert weights.

\paragraph{LLM-Enhanced Encoders.} Large language models have been applied to depression detection via instruction-tuned text encoders \citep{zhao2025multimodalllm} and audio-language models \citep{alam2025wav2vec}. We evaluate LLM-enhanced encoders in our ablation track (levels L1--L5).

\paragraph{Depression Detection on DAIC-WOZ.} The DAIC-WOZ corpus \citep{gratch2014distress} contains 189 clinical interviews with PHQ-8 labels. State-of-the-art methods include: Burdisso et al. GCN (F1=0.85) \citep{burdisso2024gcndp}; Zhang et al. MIL (AUROC=0.78) \citep{zhang2025mil}; Niu et al. multimodal fusion (F1=0.92) \citep{niu2021multimodal}; Dai et al. full fusion (F1=0.96) \citep{dai2021fullfusion}. Direct metric comparison is limited: most papers report F1 while we report AUROC.

\section{Datasets}
\label{sec:datasets}

\paragraph{DAIC-WOZ} contains 189 sessions (107 train, 35 val, 47 test) of human-agent interviews for depression assessment. Each session has a PHQ-8 binary label (depression $\ge$ 10). We use subject-independent splits: no participant appears in more than one split. Features: RoBERTa-base text (768-dim), WavLM-large audio (1536-dim), eGeMAPS (88-dim), ViT-B/16 video (1536-dim), OpenFace AUs (70-dim).

\paragraph{CMU-MOSEI} \citep{zadeh2018multimodal} contains 22,777 utterance-level samples (16,265 train, 1,869 val, 4,643 test). Labels: continuous sentiment and multi-label emotion (happiness, sadness, anger, fear, disgust, surprise). MOSEI is 120$\times$ larger than DAIC, mitigated by temperature-balanced sampling (T=2.0).

\paragraph{ChaLearn First Impressions} \citep{escalante2020chalearn} contains 10,000 video clips (6,000 train, 2,000 val, 2,000 test) with Big-Five apparent personality labels. Personality is not the same construct as clinical depression; FI serves as auxiliary supervision.

\paragraph{Leakage Protocol.} All splits are subject-independent. The graph is constructed separately for each split (split-local mode) or only on training data (inductive mode) to prevent test-label leakage through graph edges.

\section{Architecture}
\label{sec:architecture}

\begin{table}[htbp]
\centering
\caption{Notation.}
\small
\begin{tabular}{ll}
\toprule
Symbol & Meaning \\
\midrule
$K$ & Nearest neighbors (graph) \\
$M$ & Number of experts (8) \\
$d$ & Embedding dimension (256) \\
$r$ & Low-rank bottleneck (8--16) \\
$\lambda$ & Graph routing weight \\
$\sigma_t$ & Task-specific uncertainty \\
$G$ & KNN graph \\
$r_i$ & Graph routing vector \\
\bottomrule
\end{tabular}
\label{tab:notation}
\end{table}

\paragraph{Modality Encoders.} Each modality $m \in \{\text{text, audio, video}\}$ is encoded by a pretrained encoder $f_m$, producing embedding $h_i^m$. A projection layer $P_m$ maps to a common $d=256$-dimensional space:
\begin{equation}
    h_i^{m,\text{proj}} = P_m(f_m(x_i^m)) \in \mathbb{R}^{256}
\end{equation}

\paragraph{Gated Late Fusion.} Given projected modality embeddings, we apply a modality mask $m_i \in \{0,1\}^3$ and compute learned gate weights:
\begin{equation}
    \alpha_i = \text{softmax}(W_\alpha \cdot \text{concat}(h_i^{\text{text}}, h_i^{\text{audio}}, h_i^{\text{video}}) + b_\alpha)
\end{equation}
\begin{equation}
    h_i^{\text{fused}} = \sum_{m} \frac{\alpha_i \odot m_i}{\|\alpha_i \odot m_i\|_1 + \epsilon}[m] \cdot h_i^{m,\text{proj}}
\end{equation}

\paragraph{MMoEEx Expert Bank.} We instantiate $M=8$ experts $\{E_1, \ldots, E_8\}$, each an MLP with hidden dimension 256. Two experts are shared, six are task-exclusive. Expert diversity is encouraged via orthogonality regularization. Task-specific gates $g_t: \mathbb{R}^{256} \to \mathbb{R}^M$ produce expert weights.

\paragraph{KNN Graph Construction.} We construct an undirected KNN graph $G = (V, E)$ over all training samples using cosine similarity. Three construction modes are tested:
\begin{itemize}
    \item \textbf{Inductive (V0, V3)}: Graph on training only. Test nodes connect to K nearest training neighbors.
    \item \textbf{Split-local (V1, V4)}: Separate graphs per split. No cross-split edges.
    \item \textbf{Transductive (V2)}: Graph on all splits. For ablation only.
\end{itemize}

\paragraph{GraphSAGE Router.} Two-layer neighborhood aggregation produces a graph routing vector:
\begin{equation}
    h_i^{(1)} = \sigma(W^{(1)} \cdot \text{Mean}(\{h_j^{\text{fused}} : j \in \mathcal{N}(i)\} \cup \{h_i^{\text{fused}}\}))
\end{equation}
\begin{equation}
    r_i = \text{softmax}(W_r h_i^{(2)} + b_r) \in \mathbb{R}^M
\end{equation}

The final expert weights combine task gate and graph router in log-space:
\begin{equation}
    \tilde{w}_{t,i} = \text{softmax}(\log g_t(h_i) + \lambda \cdot \log r_i)
\end{equation}

\paragraph{Task Heads.} Depression: binary classifier (BCE). Sentiment: regressor (NLL). Emotion: multi-label classifier (BCE). Personality: five regression heads (NLL).

\section{Experimental Setup}
\label{sec:setup}

Training uses AdamW (lr=$1e-3$), cosine annealing, early stopping (patience=10). Temperature-balanced sampling (T=2.0) prevents MOSEI dominance. The loss combines per-task losses with learned uncertainty weights \citep{kendall2017uncertainties}:
\begin{equation}
    \mathcal{L} = \sum_t \frac{1}{2\sigma_t^2} \mathcal{L}_t + \log \sigma_t + \lambda_{\text{excl}} \mathcal{L}_{\text{excl}} + \lambda_{\text{L2}} \mathcal{L}_{\text{L2}}
\end{equation}

Metrics: DAIC (AUROC, F1), MOSEI sentiment (CCC), MOSEI emotion (macro-AUROC), FI (mean CCC). All results report 95\% bootstrap CIs (1,000 iterations).

\section{Results}
\label{sec:results}

\paragraph{Unimodal Baselines.} Text is the strongest modality for DAIC (AUROC=0.5346) and MOSEI (CCC=0.5123). Video is strongest for FI (Avg CCC=0.4578). All three datasets beat trivial baselines.

\paragraph{Fusion Ablation.} Cross-attention does not outperform gated fusion on any dataset. On DAIC, CrossAttn (AUROC=0.3117) is 0.2229 below text baseline. On MOSEI, Gated (CCC=0.6229) outperforms CrossAttn (CCC=0.5397). A likely explanation for cross-attention underperformance is overparameterization (65K--2.8M parameters vs 57K--827K for gated).

\paragraph{MMoEEx vs Baselines.} MMoEEx improves FI personality (+0.12 Avg CCC) but degrades DAIC and MOSEI sentiment, likely due to the small DAIC training set (n=107) being insufficient for joint MoE optimization.

\paragraph{Graph Routing Ablation.} Five graph variants (V0--V4) are compared in Table~\ref{tab:graph_results}.

\begin{table}[htbp]
\centering
\caption{Graph routing ablation across five variants (V0--V4). Best per-column in \textbf{bold}.}
\small
\begin{tabular}{lcccc}
\toprule
Variant & DAIC AUROC & MOSEI CCC & MOSEI AUC & FI Avg CCC \\
\midrule
V0 (inductive-k10) & 0.7124 & \textbf{0.6803} & 0.7562 & 0.4395 \\
V1 (split-local-k10) & 0.6345 & 0.5436 & 0.5467 & 0.2962 \\
V2 (transductive-k10) & 0.8505 & 0.3419 & 0.7606 & 0.3442 \\
V3 (inductive-k15) & \textbf{0.8967} & 0.5198 & 0.5985 & 0.2309 \\
V4 (split-local-k15) & 0.8351 & 0.5539 & 0.5872 & \textbf{0.5032} \\
\bottomrule
\end{tabular}
\label{tab:graph_results}
\end{table}

V0 achieves the best MOSEI sentiment (CCC=0.6803, +0.18 over MMoEEx). V3 achieves the best DAIC AUROC (0.8967, +0.40 over MMoEEx). No single variant dominates all tasks; we recommend V0 as default and V3 when DAIC is primary.

\paragraph{LLM Ablations.} LLM encoders (Mistral-7B, CLAP, LLaVA) improve over classical encoders across all tasks. L3 (Mistral + CLAP) achieves the best LLM-level DAIC AUROC (0.7210). However, no LLM level surpasses the V0 graph router, confirming that graph routing and LLM encoders provide complementary benefits.

\paragraph{Domain Adaptation.} Unsupervised domain adaptation (CORAL, MMD, DANN) produces negative transfer on clinical interview data, with plain multitask learning outperforming all DA methods.

\section{Discussion}
\label{sec:discussion}

\paragraph{What Worked.} Graph routing consistently improves over non-graph baselines, with V0 (inductive, K=10) providing the best MOSEI sentiment and V3 (inductive, K=15) the best DAIC depression detection. The KNN voting baseline confirms that GraphSAGE learned aggregation adds substantial value beyond simple neighborhood smoothing (DAIC +0.082, MOSEI +0.237). Gated late fusion provides the best fusion results on MOSEI. MMoEEx alone benefits FI personality (+0.12 CCC).

\paragraph{What Did Not.} Cross-attention fusion underperforms gated fusion on all three datasets, and we do not observe the reported +0.041 AUROC gain. MMoEEx degrades DAIC and MOSEI sentiment, likely due to insufficient training data. Domain adaptation methods produce negative transfer. Cross-demographic transfer (young to elderly) produces below-random performance, suggesting age-specific biomarkers.

\paragraph{Limitations.} DAIC's small sample size (n=107 training) limits statistical power. MOSEI audio/video features were not fully extracted. LLM ablation L6--L9 remain pending. All results are retrospective and require prospective clinical validation.

\section{Conclusion}
\label{sec:conclusion}

We presented a unified multimodal architecture with graph-gated MoE routing for mental health assessment. Graph routing consistently improves over non-graph baselines, with V0 achieving MOSEI CCC=0.6803 and V3 achieving DAIC AUROC=0.8967. Cross-attention fusion does not replicate a reported gain in our setting. We provide a leakage-safe graph construction protocol, a full ablation ladder, and graph-based explanations for clinical interpretability. Code and artifacts are released for reproducibility.

\paragraph{Ethics Statement.} This system is intended as a decision-support tool only and should not replace clinical diagnosis. All results require prospective validation before clinical deployment. Apparent personality and clinical depression are distinct constructs and should not be conflated.

\bibliographystyle{unsrtnat}
\bibliography{references/bibliography}

\end{document}
